\subsection{Сценарий 4: Simple Tag}

Как уже упоминалось \hyperref[exp-st]{выше}, этот сценарий не подразумевает коммуникационных действий. Но он был выбран нами, так как он достаточно сложный, и при этом агенты из одной команды имеют одинаковые пространства наблюдения и действий, что позволяет применить вариант алгоритма MADDPG с общим мозгом.

Мы поставили эксперимент с двумя преследователями и одной жертвой.

\begin{figure}[ht!]
	\center
%	\includegraphics [scale=0.6] {my_folder/images/ch5/st-rew.png}
	\includegraphics [scale=0.6] {my_folder/images/in_progress.png}
	\caption{Среднее вознаграждение преследователей и жертвы с алгоритмами DDPG, MADDPG и MADDPG с общим мозгом.}
	\label{fig:result-st-rew}
\end{figure}

\begin{figure}[ht!]
	\center
%	\includegraphics [scale=0.6] {my_folder/images/ch5/st-4vs2-rew.png}
	\includegraphics [scale=0.6] {my_folder/images/in_progress.png}
	\caption{График консистентности действий коммуникации для алгоритмов DDPG, MADDPG и MADDPG с общим мозгом.}
	\label{fig:result-st-4vs2-rew}
\end{figure}

\begin{table}[t!]
	\centering\small
	\caption{Среднее время, потраченное на обучение с различными алгоритмами в 20000 эпизодов.}
	\label{tab-sr-time}
	\begin{tabular}{|l|l|l|l|l|l|}
		\hline
		&DDPG&MADDPG&MADDPG с одним мозгом\\
		\hline
		Время, 20000 эпизодов&7ч. 46мин.&10ч. 18мин.&8ч. 20мин.\\ \hline
	\end{tabular}
	\normalsize% возвращаем шрифт к нормальному
\end{table}